\chapter{SYSTEM DESIGN AND HARDWARE IMPLEMENTATION}\label{ch:system}

\section{Overall Architecture}

Figure~\ref{fig:arch} shows the final system (kit v6). Raw signals arrive second by second. On the AMD Kria KV260 the
Arm cores only move the raw samples into the FPGA; the feature engine and the classifier both run in the fabric. When
the alarm turns on, the FPGA drives a red LED and a buzzer on a breadboard, and an ESP32 on the same breadboard sends
a push notification to the doctor's phone. A laptop browser shows every stage of the detector live.

\begin{figure}[H]
\centering
\resizebox{\textwidth}{!}{% SeizeLUT on the KV260 (kit v6): processing system (Arm) and programmable logic (fabric)
\begin{tikzpicture}[font=\scriptsize, >=stealth,
  blk/.style={draw, rounded corners=1.5pt, align=center, minimum height=5.2mm, inner sep=1.6pt, fill=white},
  grp/.style={draw, dashed, rounded corners=3pt, inner sep=3pt}]
  % Arm side
  \node[blk] (sens) at (0,0) {bte-EEG$\times$2, ECG,\\EMG, ACC, GYR};
  \node[blk, right=3mm of sens] (drv) {driver: raw samples\\$\to$ AXI4-Lite};
  \node[grp, fit=(drv), label={[font=\scriptsize\bfseries]above:Arm Cortex-A53 (PS)}] (ps) {};
  % fabric
  \node[blk, below=8mm of sens.south west, anchor=north west] (fe) {feature engine\\shift-add, 2,323 instr.\\38 features / s};
  \node[blk, right=2.5mm of fe] (cmp) {146\\comparators};
  \node[blk, right=2.5mm of cmp] (lut) {50 LUT6\\(DWN)};
  \node[blk, right=2.5mm of lut] (pop) {popcount\\score};
  \node[blk, right=2.5mm of pop] (alm) {decision,\\$k$-of-10 alarm};
  \node[blk, below=3.5mm of cmp] (rom) {BIST ROM\\255 + 769 vectors};
  \node[blk, below=3.5mm of lut, xshift=5mm] (fi) {fault injection\\(stuck-at, TT flip)};
  \node[blk, below=3.5mm of alm] (chk) {BIST check: score,\\decision, alarm};
  \node[blk, below=3.5mm of fe] (io) {alarm out: LEDs,\\buzzer, button};
  \node[grp, fit=(fe)(alm)(rom)(io)(chk), label={[font=\scriptsize\bfseries]below:programmable logic, 100 MHz, no multiplier}] (pl) {};
  \draw[->] (sens) -- (drv);
  \draw[->] (drv.south) |- ++(0,-3mm) -| (fe.north);
  \draw[->] (fe) -- (cmp);
  \draw[->] (cmp) -- (lut);
  \draw[->] (lut) -- (pop);
  \draw[->] (pop) -- (alm);
  \draw[->] (rom.north) -- (cmp.south);
  \draw[->] (fi.north -| lut.south) -- (lut.south);
  \draw[->] (alm.south) -- (chk.north);
\end{tikzpicture}
}
\caption{Block diagram of SeizeLUT on the Kria KV260, with the fault-injection registers and the self-test that let
the device check that the fabric still computes the validated classifier}\label{fig:arch}
\end{figure}

\section{Hardware Components}

\begin{table}[H]
\centering\small
\caption{Main hardware components}\label{tab:components}
\begin{tabularx}{\textwidth}{|P{30mm}|P{50mm}|L|}
\hline
\multicolumn{1}{|c|}{\textbf{Component}} & \multicolumn{1}{c|}{\textbf{Part / specification}} & \multicolumn{1}{c|}{\textbf{Function}}\\
\hline
FPGA board & AMD Kria KV260 (Zynq UltraScale+ XCK26: four Arm Cortex-A53 cores and FPGA fabric), Ubuntu 22.04 with Kria-PYNQ & Feature engine, classifier, self-test, fault injection, demo outputs\\
\hline
Alarm LED & Red LED with 220\,$\Omega$ resistor on PMOD B1 & Lights while the alarm is on\\
\hline
Buzzer & Buzzer switched by a BC547 transistor (1\,k$\Omega$ base resistor) on PMOD B2 & Sounds while the alarm is on\\
\hline
Alert module & ESP32 Dev Module, USB serial link to the board, alarm wire from PMOD B1 to GPIO\,27 & Push notification and live page on the doctor's phone\\
\hline
Phone & iPhone with the ntfy app; personal hotspot as the network & Receives the alert; ``On my way'' reply\\
\hline
Laptop & Web browser & Live dashboard of the detector\\
\hline
\end{tabularx}
\end{table}

The board runs the detector from recorded test-patient signals stored on it, second by second, exactly as it would
from live sensors. The ESP32 is powered by the USB cable from the board, which also carries the alarm messages; the
alarm wire is a backup that works even if the board's software stops.

\section{Sensor Configuration}

The detector uses the signals of the SeizeIT2 wearable (Table~\ref{tab:modalities}). The two EEG channels are one
behind-the-ear channel and one cross-head channel, chosen per recording by the dataset. A presence bit per sensor
tells the classifier when a sensor is missing.

\begin{table}[H]
\centering\small
\caption{Signals recorded by the wearable (rates measured from the files)}\label{tab:modalities}
\begin{tabular}{|l|l|c|}
\hline
\multicolumn{1}{|c|}{\textbf{Signal}} & \multicolumn{1}{c|}{\textbf{Channels}} & \textbf{Rate}\\
\hline
Behind-the-ear EEG & 2 (left, right or cross-head) & 256 Hz\\
\hline
ECG & 1 & 256 Hz\\
\hline
EMG & 1 & 256 Hz\\
\hline
Motion & 3-axis accelerometer + 3-axis gyroscope & 25 Hz\\
\hline
\end{tabular}
\end{table}

\section{Board Software and Live Demonstration}

\begin{itemize}
\item \textbf{Driver and tests:} a Python driver (PYNQ) with a fast C path writes raw samples and reads results. One
  script (\texttt{run\_board\_tests.sh}) runs nine tests and saves every number to a results file.
\item \textbf{Dashboard:} recorded test-patient signals stream second by second through the board, and a laptop
  browser shows the signals, features, comparator bits, LUT votes, score and alarm, with a check on every window that
  the FPGA equals the software (Figure~\ref{fig:dashboard}, Chapter~\ref{ch:results}).
\item \textbf{Breadboard:} the alarm LED, the buzzer and the ESP32 alarm wire are driven directly by the FPGA pins,
  so they follow the hardware decision.
\end{itemize}

\begin{table}[H]
\centering\small
\caption{Breadboard wiring of the demonstration}\label{tab:wiring}
\begin{tabularx}{\textwidth}{|P{30mm}|L|}
\hline
\multicolumn{1}{|c|}{\textbf{PMOD pin}} & \multicolumn{1}{c|}{\textbf{Connection}}\\
\hline
B1 & 220\,$\Omega$ to the red LED (long leg); short leg to ground; also ESP32 GPIO\,27\\
\hline
B2 & 1\,k$\Omega$ to the BC547 base; emitter to ground; collector to the buzzer ($-$)\\
\hline
B3 (optional) & ESP32 GPIO\,26 (second alarm wire)\\
\hline
3.3\,V & buzzer (+)\\
\hline
GND & ground rail, ESP32 GND\\
\hline
\end{tabularx}
\end{table}

\section{Doctor Alert with an ESP32}\label{sec:alert}

\begin{itemize}
\item \textbf{Board to ESP32:} the board sends one short message per event over USB serial (alarm start, updates
  every 15\,s, escalations, end, patient button).
\item \textbf{ESP32 to phone:} the ESP32 joins the phone's Wi-Fi hotspot, sends a push notification through the ntfy
  service and serves a live alarm page on the same network.
\item \textbf{Phone to board:} the doctor's ``On my way'' reaches the ESP32 and is sent back to the board, where the
  dashboard shows it.
\end{itemize}

The detector answers ``seizure-like or not''; it was not trained to name a seizure type, so the alert reports only
signs that the front end measures, compared with the minute before the alarm (Table~\ref{tab:signs}).

\begin{table}[H]
\centering\small
\caption{The signs in the alert}\label{tab:signs}
\begin{tabularx}{\textwidth}{|l|L|P{45mm}|}
\hline
\multicolumn{1}{|c|}{\textbf{Sign}} & \multicolumn{1}{c|}{\textbf{Shown as}} & \multicolumn{1}{c|}{\textbf{Measured from}}\\
\hline
Heart rate & before and now, NORMAL / HIGH / VERY HIGH & ECG R-peak intervals\\
\hline
Shaking & LITTLE / SOME / STRONG & accelerometer (or gyroscope)\\
\hline
Muscle tension & NORMAL / HIGH & EMG power\\
\hline
Brain activity & MILD / HIGH, slow or fast waves & EEG line length and band power\\
\hline
Pattern hint & ``motor pattern likely'' or ``little movement'' & movement change\\
\hline
\end{tabularx}
\end{table}

Because the doctor has very little time, the alert is designed to be read in about two seconds:

\begin{itemize}
\item \textbf{Notification:} three lines, for example ``SEIZURE Bed 3 / GO TO BED 3 NOW: PROTECT FROM INJURY / Heart
  120 bpm HIGH \textbar{} Shaking STRONG \textbar{} Brain HIGH''.
\item \textbf{Page:} status and a running timer, one instruction, ``On my way'' and ``Call'' buttons, and four colour
  tiles (red = act, amber = watch, green = normal); the emergency plan and first-aid steps are folded below.
\item \textbf{One instruction}, the most urgent first: CHECK BREATHING AND PULSE NOW (heart rate above 150 or below
  40\,bpm, or an unreadable ECG without shaking) $>$ OVER 5 MIN: START EMERGENCY PLAN $>$ OVER 2 MIN: GET EMERGENCY
  PLAN READY $>$ GO TO BED NOW: PROTECT FROM INJURY (strong shaking) $>$ GO TO BED NOW. When it ends: SEIZURE STOPPED:
  CHECK BREATHING, TURN ON SIDE.
\item \textbf{Escalation} at 2 and 5 minutes of alarm, then every minute; five minutes is the ILAE time point for
  treating a convulsive seizure as status epilepticus. \textbf{Reminders} every minute until someone taps ``On my
  way''; \textbf{clusters} (``2nd seizure this hour'') are flagged; only a bed name is sent, never a patient name.
\end{itemize}

The movement-based pattern hint was checked against the doctors' motor and non-motor labels of the 702 seizures of
the training patients only. In the first 30\,s movement rose 4.8-fold (median) in motor seizures and 1.04-fold in
non-motor seizures. A threshold (2.6-fold) chosen on patients 1--48 matched the doctors' label on patients 49--96 in
64\,\% of cases (balanced accuracy; 48\,\% of motor and 79\,\% of non-motor seizures), so the hint is useful but is not
a diagnosis.
